% LaTeX supplement: Additional sections for dissertation

\section{Implementation Architecture}
\label{sec:implementation}

\subsection{System Architecture Overview}

The LLM-based refactoring system implements a modular, pipeline-oriented architecture consisting of five primary subsystems:

\begin{enumerate}
    \item \textbf{Smell Detection Layer}: Integrates Steel and SNutsJS detectors
    \item \textbf{LLM Orchestration Layer}: Multi-provider client with strategy pattern
    \item \textbf{Database Layer}: SQLite-based experiment repository
    \item \textbf{Execution Engine}: Automated test runner and validation framework
    \item \textbf{Analysis Layer}: Statistical comparison and metric computation
\end{enumerate}

\subsection{Multi-Provider LLM Client Architecture}

The system employs a factory pattern to abstract provider-specific API differences:

\begin{lstlisting}[language=Python, caption=LLM Client Factory Pattern]
class BaseLLMClient(ABC):
    """Abstract interface for all providers"""
    @abstractmethod
    def generate(self, model, system_prompt, user_prompt,
                 temperature, top_p, max_tokens) -> Dict:
        pass

class OpenAICompatibleClient(BaseLLMClient):
    """Supports: HuggingFace, OpenAI, Together, Fireworks"""
    def generate(self, ...):
        client = OpenAI(api_key=self.api_key, 
                       base_url=self.base_url)
        # GPT-5.x uses Responses API
        if model.startswith("gpt-5"):
            response = client.responses.create(...)
        else:
            response = client.chat.completions.create(...)
        return {"content": ..., "tokens": ..., "latency": ...}

class AnthropicClient(BaseLLMClient):
    """Native Anthropic SDK for Claude models"""
    def generate(self, ...):
        client = anthropic.Anthropic(api_key=self.api_key)
        response = client.messages.create(...)
        return {"content": ..., "tokens": ..., "latency": ...}

class LLMClientFactory:
    @staticmethod
    def create(provider: str, api_key: str, 
               base_url: str) -> BaseLLMClient:
        if provider == LLMProvider.ANTHROPIC:
            return AnthropicClient(api_key, base_url)
        else:
            return OpenAICompatibleClient(api_key, base_url)
\end{lstlisting}

\textbf{Design Rationale:}
\begin{itemize}
    \item Isolates provider-specific authentication, request formatting, and response parsing
    \item Enables seamless addition of new providers without modifying core refactoring logic
    \item Standardizes token counting and latency measurement across providers
\end{itemize}

\subsection{Prompt Construction Pipeline}

Dynamic prompt generation follows a template-based approach with strategy-specific augmentation:

\begin{algorithm}
\caption{Prompt Construction}
\label{alg:prompt-construction}
\begin{algorithmic}[1]
\Require $smell\_context$, $strategy\_id$
\Ensure $prompt = \{system, user\}$
\State $base \gets \text{LoadTemplate}(strategy\_id)$
\State $definition \gets smell\_context.\text{definition}$
\State $code \gets smell\_context.\text{code\_snippet}$
\If{$strategy\_id = \text{ZERO\_SHOT}$}
    \State $user \gets base + definition + code$
\ElsIf{$strategy\_id = \text{FEW\_SHOT}$}
    \State $examples \gets smell\_context.\text{examples}[1:3]$
    \State $user \gets base + definition + \text{FormatExamples}(examples) + code$
\ElsIf{$strategy\_id = \text{COT}$}
    \State $detection \gets smell\_context.\text{detection\_criteria}$
    \State $strategies \gets smell\_context.\text{refactoring\_strategies}$
    \State $examples \gets smell\_context.\text{examples}[1:2]$
    \State $cot\_template \gets \text{LoadCoTReasoning}()$
    \State $user \gets base + definition + detection + strategies + examples + code + cot\_template$
\EndIf
\State $system \gets \text{"You are a JavaScript test refactoring expert..."}$
\State \Return $\{system, user\}$
\end{algorithmic}
\end{algorithm}

\subsection{Code Extraction and Validation}

LLM responses undergo multi-stage parsing to extract executable code:

\begin{lstlisting}[language=Python, caption=Code Extraction Pipeline]
def extract_code_from_response(text: str) -> str:
    """Extract JavaScript code from LLM response"""
    # Stage 1: Markdown fence extraction
    code_blocks = re.findall(
        r'```(?:javascript|js)?\n(.*?)```', 
        text, 
        re.DOTALL | re.IGNORECASE
    )
    
    if code_blocks:
        # Prefer explicitly marked JavaScript blocks
        return clean_extracted_code(code_blocks[0])
    
    # Stage 2: Fallback - entire response if no fences
    clean_text = re.sub(r'^[#\s]*(?:javascript|js)\s*', '', text)
    return clean_extracted_code(clean_text)

def clean_extracted_code(code: str) -> str:
    """Remove LLM artifacts and formatting"""
    # Remove leading/trailing markers
    code = re.sub(r'^```(?:javascript|js)?\s*', '', code)
    code = re.sub(r'```\s*$', '', code)
    # Remove explanatory comments added by LLM
    code = re.sub(r'^\s*//\s*(Refactored|Complete).*$', '', 
                  code, flags=re.MULTILINE)
    return code.strip()
\end{lstlisting}

\subsection{Backup and Restoration System}

To ensure repository integrity during experiments, the system implements a transactional backup mechanism:

\begin{lstlisting}[language=Python, caption=Backup Manager]
class BackupManager:
    def replace_snippet(self, repo_name, file_path, 
                       original_snippet, refactored_snippet):
        """Atomic snippet replacement with backup"""
        # 1. Create timestamped backup
        backup_path = self._create_backup(repo_name, file_path)
        
        # 2. Validate snippet exists in file
        full_path = self.repos_dir / repo_name / file_path
        content = full_path.read_text()
        if original_snippet not in content:
            raise SnippetReplacementError("Original not found")
        
        # 3. Replace and write
        new_content = content.replace(
            original_snippet, 
            refactored_snippet, 
            1  # Replace only first occurrence
        )
        full_path.write_text(new_content)
        
        return backup_path
    
    def undo_refactor(self, repo_name, file_path):
        """Restore from most recent backup"""
        backup = self._find_latest_backup(repo_name, file_path)
        shutil.copy(backup, self.repos_dir / repo_name / file_path)
\end{lstlisting}

\textbf{Backup Structure:}
\begin{verbatim}
.backup/
├── repository_name/
│   ├── path/to/test_file.test.js.20260222_143052.bak
│   ├── path/to/test_file.test.js.20260222_143053.bak
│   └── ...
\end{verbatim}

\subsection{Smell Analysis Framework}

Post-execution analysis compares baseline and refactored smell detection results:

\begin{lstlisting}[language=Python, caption=Smell Analysis]
def compare_repositories(baseline_df, refactored_df, 
                        target_file, target_smell):
    """Compare smell detection results"""
    # Filter to target file
    baseline_smells = baseline_df[
        baseline_df['file'] == target_file
    ]
    refactored_smells = refactored_df[
        refactored_df['file'] == target_file
    ]
    
    # Target smell removal check
    target_baseline = baseline_smells[
        baseline_smells['smell_type'] == target_smell
    ]
    target_refactored = refactored_smells[
        refactored_smells['smell_type'] == target_smell
    ]
    
    target_removed = (
        len(target_baseline) > 0 and 
        len(target_refactored) == 0
    )
    
    # New smell introduction check
    baseline_types = set(baseline_smells['smell_type'])
    refactored_types = set(refactored_smells['smell_type'])
    
    introduced_smells = refactored_types - baseline_types
    introduced_new = len(introduced_smells) > 0
    
    return {
        'target_smell_removed': target_removed,
        'introduced_new_smells': introduced_new,
        'new_smell_types': list(introduced_smells),
        'baseline_count': len(baseline_smells),
        'refactored_count': len(refactored_smells)
    }
\end{lstlisting}

\subsection{Test Execution and Coverage Analysis}

The system executes test suites and parses Jest output for comprehensive metrics:

\begin{lstlisting}[language=Python, caption=Test Results Parsing]
def extract_coverage_summary(output: str) -> str:
    """Parse Jest coverage table"""
    match = re.search(
        r'(-+\|){4,}.*?(-+\|){4,}', 
        output, 
        re.DOTALL
    )
    return match.group(0) if match else None

def extract_test_results(output: str) -> str:
    """Parse test execution summary"""
    patterns = [
        r'Tests:\s+\d+\s+passed.*',
        r'Test Suites:\s+\d+\s+passed.*',
        r'Time:\s+[\d.]+s'
    ]
    
    results = []
    for pattern in patterns:
        match = re.search(pattern, output)
        if match:
            results.append(match.group(0))
    
    return '\n'.join(results)

def parse_test_counts_from_summary(summary: str) -> Dict:
    """Extract numeric test metrics"""
    test_match = re.search(
        r'Tests:\s+(\d+)\s+passed,\s+(\d+)\s+total',
        summary
    )
    suite_match = re.search(
        r'Test Suites:\s+(\d+)\s+passed.*?(\d+)\s+total',
        summary
    )
    
    return {
        'tests_passed': int(test_match.group(1)),
        'tests_total': int(test_match.group(2)),
        'test_suites_passed': int(suite_match.group(1)),
        'test_suites_total': int(suite_match.group(2))
    }
\end{lstlisting}

\subsection{Database Schema Details}

\begin{table}[h]
\centering
\caption{Complete Experiment Table Schema}
\small
\begin{tabular}{@{}llp{5cm}@{}}
\toprule
\textbf{Column} & \textbf{Type} & \textbf{Description} \\ \midrule
id & INTEGER & Primary key \\
study\_smell\_id & INTEGER & FK to study\_smells \\
file\_id & INTEGER & FK to files \\
ai\_tool & STRING & Provider name (HuggingFace/OpenAI/Anthropic) \\
ai\_model\_version & STRING & Model display name \\
prompting\_approach & STRING & Strategy (Zero-Shot/Few-Shot/CoT) \\
original\_code & TEXT & Baseline smelly code \\
refactored\_code & TEXT & LLM-generated code \\
prompt\_text & TEXT & Full prompt sent (optional) \\
refactoring\_completed & BOOLEAN & Phase 1 success \\
smell\_removed & BOOLEAN & Target smell eliminated \\
introduced\_new\_smells & BOOLEAN & New smells detected \\
tests\_still\_passing & BOOLEAN & Test suite status \\
coverage\_decreased & BOOLEAN & Coverage regression \\
tests\_changed & BOOLEAN & Test count/status changed \\
tests\_pass\_rate\_decreased & BOOLEAN & Pass rate regression \\
refactor\_phase\_completed & BOOLEAN & Phase 1 checkpoint \\
execution\_phase\_completed & BOOLEAN & Phase 2 checkpoint \\
tokens\_used & INTEGER & Total LLM tokens \\
llm\_latency\_seconds & FLOAT & API response time \\
execution\_time\_seconds & FLOAT & Total experiment time \\
created\_at & DATETIME & Record creation \\
updated\_at & DATETIME & Last modification \\ \bottomrule
\end{tabular}
\end{table}

\subsection{Batch Experiment Orchestration}

For large-scale experiments, the system supports batch processing with progress tracking:

\begin{lstlisting}[language=Python, caption=Batch Execution]
def run_batch_experiments(strategy_id, model_id, limit=None):
    """Execute refactoring experiments in batch"""
    smells = get_pending_smells(strategy_id, model_id, limit)
    
    stats = {'total': len(smells), 'completed': 0, 'failed': 0}
    failed_experiments = []
    
    for idx, (smell_id, repo, file_path, smell_type) in enumerate(smells):
        try:
            # Phase 1: Refactor
            result = refactor_smell(smell_id, strategy_id, model_id)
            experiment_id = result['experiment_id']
            
            # Phase 2: Execute
            validation = execute_experiment(experiment_id)
            
            stats['completed'] += 1
            
        except Exception as e:
            stats['failed'] += 1
            failed_experiments.append((smell_id, str(e)))
            
        # Progress reporting
        progress = (idx + 1) / stats['total'] * 100
        print(f"Progress: {progress:.1f}% ({idx+1}/{stats['total']})")
    
    # Generate summary report
    write_batch_summary(strategy_id, model_id, stats, 
                       failed_experiments)
\end{lstlisting}

\subsection{Performance Optimizations}

Several optimizations ensure scalability:

\begin{enumerate}
    \item \textbf{Parallel Smell Detection}: Steel and SNutsJS run concurrently via \texttt{asyncio}
    \item \textbf{Database Caching}: Baseline smell results cached to avoid redundant queries
    \item \textbf{Incremental Commits}: Results committed per experiment to enable interruption/resumption
    \item \textbf{Manifest-Based Execution}: Phase 2 can resume from Phase 1 manifest files
    \item \textbf{Connection Pooling}: Database session management with proper cleanup
\end{enumerate}

\begin{lstlisting}[language=Python, caption=Async Smell Detection]
async def run_detectors_parallel(repo_path, output_dir):
    """Run smell detectors concurrently"""
    async def run_snuts_async():
        # Execute SNutsJS detector
        result = await asyncio.to_thread(
            run_snuts_detector, repo_path, output_dir
        )
        return ('snuts', result['success'], result['message'])
    
    async def run_steel_async():
        # Execute Steel detector
        result = await asyncio.to_thread(
            run_steel_detector, repo_path, output_dir
        )
        return ('steel', result['success'], result['message'])
    
    # Concurrent execution
    results = await asyncio.gather(
        run_snuts_async(),
        run_steel_async(),
        return_exceptions=True
    )
    
    return results
\end{lstlisting}

\subsection{Error Handling and Recovery}

The system implements comprehensive error handling:

\begin{itemize}
    \item \textbf{Network Errors}: Automatic retry with exponential backoff for transient API failures
    \item \textbf{Extraction Failures}: Fallback parsers when primary code extraction fails
    \item \textbf{Test Timeouts}: Configurable timeout (default: 300s) with graceful termination
    \item \textbf{File Restoration}: Automatic backup restoration on any exception during Phase 2
    \item \textbf{Database Rollback}: Transaction-based updates with rollback on constraint violations
\end{itemize}

\subsection{Re-execution Support (--redo Flag)}

To enable experiment re-runs (e.g., after fixing issues), the system supports execution phase reset:

\begin{lstlisting}[language=Python, caption=Experiment Reset]
def reset_experiment_execution_data(session, experiment_id):
    """Clean execution phase data for re-run"""
    # Delete all test_results records
    session.query(TestResult).filter_by(
        experiment_id=experiment_id
    ).delete()
    
    # Reset analysis flags
    update_experiment(
        session, experiment_id,
        execution_phase_completed=False,
        tests_still_passing=None,
        smell_removed=None,
        introduced_new_smells=None,
        coverage_decreased=None,
        tests_changed=None,
        tests_pass_rate_decreased=None
    )
    
    session.flush()
\end{lstlisting}

\textbf{Preservation Guarantees:}
\begin{itemize}
    \item Refactored code preserved (Phase 1 data intact)
    \item Experiment ID and associations maintained
    \item LLM metrics (tokens, latency) retained
    \item Only execution-phase results cleared
\end{itemize}

This ensures experiments can be re-validated without re-calling LLM APIs, supporting:
\begin{enumerate}
    \item Test suite updates (new tests added to repository)
    \item Smell detector improvements (re-run detection with updated tools)
    \item Analysis bug fixes (correct metric calculation errors)
\end{enumerate}

\section{Research Implications}

\subsection{Contributions to SE Literature}

This methodology advances software engineering research in three dimensions:

\begin{enumerate}
    \item \textbf{Automated Test Maintenance}: Demonstrates LLM viability for systematic test quality improvement
    \item \textbf{Prompt Engineering for Code}: Quantifies impact of strategy selection on refactoring success
    \item \textbf{LLM Comparative Analysis}: Provides empirical evidence of model-specific strengths in code transformation
\end{enumerate}

\subsection{Practical Applications}

Results inform development of:
\begin{itemize}
    \item IDE-integrated refactoring assistants
    \item CI/CD quality gates for test smell detection + automated fixes
    \item Educational tools demonstrating test best practices
\end{itemize}

\subsection{Future Research Directions}

Potential extensions include:
\begin{enumerate}
    \item \textbf{Multi-Smell Refactoring}: Handling tests with multiple concurrent smells
    \item \textbf{Cross-Language Transfer}: Evaluating methodology on Python, Java, C\# codebases
    \item \textbf{Human-in-the-Loop}: Hybrid approaches combining LLM suggestions with developer review
    \item \textbf{Fine-Tuning}: Training specialized models on curated smell-refactoring pairs
    \item \textbf{Incremental Refactoring}: Iterative improvement cycles with developer feedback
\end{enumerate}
